\documentclass[aspectratio=169,11pt]{beamer}

\usetheme{Madrid}
\usecolortheme{default}
\usefonttheme{professionalfonts}
\setbeamertemplate{navigation symbols}{}
\setbeamertemplate{footline}[frame number]

\usepackage{amsmath, amssymb, booktabs}

\title{Synthesis, Frontier Questions}
\subtitle{And Student Research Designs}
\author{Special Topic 4: Urban and Labor -- Week 6}
\date{}

\begin{document}

\begin{frame}
  \titlepage
\end{frame}

\begin{frame}{Central question}
\begin{itemize}
    \item What does a credible urban-labor research design look like?
    \item Urban labor research is distinct when it starts from work, wages, access, mobility, allocation, inequality, or welfare.
    \item The city is not the estimand. The spatial labor mechanism is.
    \item Week 6 turns the course into student research memos.
\end{itemize}
\end{frame}

\begin{frame}{Course synthesis}
\small
\begin{tabular}{p{0.23\linewidth} p{0.57\linewidth}}
\toprule
Mechanism & Labor question \\
\midrule
Commuting and access & Which jobs can workers search for, reach, accept, and keep? \\
Housing and rents & Who can enter opportunity, and who captures local gains? \\
Segregation and exposure & How do neighborhoods shape access, networks, schools, and mobility? \\
Safety and environment & How do risk and exposure change feasible work and welfare? \\
Migration and demand & How do shocks reallocate workers, firms, rents, and commuting flows? \\
\bottomrule
\end{tabular}
\end{frame}

\begin{frame}{Urban-labor architecture}
\[
Y_{igt}
=
f\left(M_{gt}, F_{igt}, I_{gt}, E_{gt}\right)
\]
\small
\begin{tabular}{p{0.18\linewidth} p{0.65\linewidth}}
\toprule
Object & Interpretation \\
\midrule
$Y_{igt}$ & labor outcome: employment, wages, access, commute, retention, mobility, welfare \\
$M_{gt}$ & urban mechanism: access, housing, exposure, safety, environment, demand \\
$F_{igt}$ & worker or firm friction: moving cost, search, wealth, networks, bargaining \\
$I_{gt}$ & institutions: zoning, transit, schools, policy, enforcement, infrastructure \\
$E_{gt}$ & equilibrium response: rents, migration, commuting, entry, exit, sorting \\
\bottomrule
\end{tabular}
\end{frame}

\begin{frame}{Project-design template}
\small
\begin{tabular}{p{0.25\linewidth} p{0.55\linewidth}}
\toprule
Design element & Guiding question \\
\midrule
Labor outcome & What is the labor object? \\
Urban mechanism & What spatial force is doing the work? \\
Unit of geography & Which geography matches the mechanism? \\
Comparison & What changes relative to what counterfactual? \\
Adjustment margins & Who responds: workers, firms, commuters, landlords, migrants? \\
Equilibrium issue & Do rents, commuting, migration, or firm entry change the estimate? \\
Welfare object & What improves: nominal pay, real access, safety, opportunity, surplus? \\
External relevance & Why should readers outside this city care? \\
\bottomrule
\end{tabular}
\end{frame}

\begin{frame}{Where the field is now}
\begin{itemize}
    \item Wages, rents, amenities, and mobility are interpreted together.
    \item Local shocks are spatial-incidence problems, not one wage coefficient.
    \item Housing supply and moving costs govern access to productive places.
    \item Access is richer than distance: transit, networks, schools, safety, stigma, and schedules matter.
    \item Welfare increasingly means real access, risk, child opportunity, and incidence.
\end{itemize}
\end{frame}

\begin{frame}{Where the field is going}
\small
\begin{tabular}{p{0.27\linewidth} p{0.50\linewidth}}
\toprule
Area & Frontier direction \\
\midrule
Spatial equilibrium & dynamic incidence under remote work, climate, and zoning reform \\
Housing and access & distributional gains from reform and who captures surplus \\
Neighborhood exposure & mechanism separation: schools, networks, safety, stigma \\
Risk and feasible work & hidden harms, exposure data, and incomplete compensation \\
Local shocks & slow adjustment, place-based policy, mobility barriers \\
Labor power & commuting-linked monopsony and spatial outside options \\
\bottomrule
\end{tabular}
\end{frame}

\begin{frame}{Frontier opportunities}
\small
\begin{tabular}{p{0.28\linewidth} p{0.50\linewidth}}
\toprule
Topic & Sample labor question \\
\midrule
Remote work & Who gains outside options and who loses local networks or service demand? \\
Housing reform & Do constrained workers gain real access to high-wage markets? \\
Climate and heat & Which workers lose productivity, hours, or mobility under exposure? \\
Commuting and power & Do travel costs raise employer wage-setting power? \\
Place-based policy & Do gains accrue to incumbents, migrants, firms, or landlords? \\
Child opportunity & Which urban changes produce adult labor-market gains? \\
Redevelopment & Are jobs created for incumbents or mainly for new residents and commuters? \\
Digital infrastructure & Does digital access broaden or concentrate urban opportunity? \\
\bottomrule
\end{tabular}
\end{frame}

\begin{frame}{From phenomenon to labor paper}
\begin{itemize}
    \item Weak: rents rose after redevelopment.
    \item Stronger: did redevelopment change incumbent workers' employment, commute burden, real access, or displacement risk?
    \item Weak: remote work changed cities.
    \item Stronger: did remote work change worker outside options by occupation, home quality, broadband, and local service dependence?
    \item Good projects convert urban change into a labor outcome and a mechanism.
\end{itemize}
\end{frame}

\begin{frame}{What makes a setting travel?}
\small
\begin{tabular}{p{0.26\linewidth} p{0.50\linewidth}}
\toprule
Weak claim & Stronger claim \\
\midrule
This city is interesting & This setting reveals a general spatial labor mechanism \\
Neighborhoods matter & The design separates transit, schools, rents, networks, safety \\
Treated workers gain & Incidence traces workers, firms, landlords, commuters, migrants \\
Wages increased & Real access, commute burden, risk, and child opportunity are assessed \\
Policy worked here & The mechanism is linked to transportability conditions \\
\bottomrule
\end{tabular}
\end{frame}

\begin{frame}{Common empirical failure modes}
\begin{itemize}
    \item Mechanism slippage: one treatment changes many channels.
    \item Geography mismatch: the unit does not match the mechanism.
    \item Ignoring equilibrium: rents, migration, commuting, and firm entry move.
    \item Nominal-outcome bias: wages stand in for welfare too quickly.
    \item Localism without mechanism: the paper documents a place, not a portable labor question.
    \item Sorting versus treatment confusion: residence and exposure are treated as random.
\end{itemize}
\end{frame}

\begin{frame}{One-page research memo}
\small
\begin{tabular}{p{0.30\linewidth} p{0.48\linewidth}}
\toprule
Line item & What it must say \\
\midrule
Question & one sentence with labor outcome and mechanism \\
Outcome & measurable labor object \\
Mechanism & spatial channel and affected margin \\
Geography & unit and why it fits \\
Comparison & source of identifying variation \\
Incidence & affected population and likely adjustment margins \\
Threats & main confounds, sorting, spillovers, measurement gaps \\
Welfare & what the estimate can and cannot say \\
Portability & why the mechanism matters elsewhere \\
\bottomrule
\end{tabular}
\end{frame}

\begin{frame}{Week 6 lab}
\begin{itemize}
    \item Reproduce: inspect synthetic examples of design architecture.
    \item Diagnose: classify common failure modes and repairs.
    \item Transfer: score frontier project ideas for labor focus, mechanism clarity, design plausibility, and portability.
    \item Write: draft a one-page urban-labor research memo.
    \item Goal: leave with a project-shaped object, not just a topic.
\end{itemize}
\end{frame}

\begin{frame}{Takeaways}
\begin{itemize}
    \item Urban labor research starts from labor outcomes and treats space as a mechanism.
    \item Strong designs match outcome, mechanism, geography, comparison, incidence, and welfare.
    \item The frontier is broad: remote work, housing reform, climate, commuting power, place policy, child opportunity, redevelopment, and digital access.
    \item Student projects should make one city reveal a mechanism that travels.
\end{itemize}
\end{frame}

\end{document}
